% TOP_PROBLEM: {"id":30007584,"problem_number":"LOCAL-30007584","title":"Biases under high stakes and experience","category_id":21,"category":{"id":21,"name":"economics","display_name":"Economics","description":"Open economic theory statements, published research agendas, and proposed scoped specifications; question provenance and historical priority are qualified per record.","slug":"economics","order_index":21,"created_at":"2026-10-08T00:00:00Z"},"status":"provenance_pending","external_url":null,"legacy_ids":[],"published":false,"statement":"Estimate how stakes and repeated experience change the specified bias measures, distinguishing learning from selection and preference changes.","economics":{"original_id":73,"field":"Behavioral and experimental economics","kind":"empirical","formalization_basis":"proposed_specification","source_status":"not_verified","priority_status":"not_established","priority_note":"The paper tests an older high-stakes question and cites Thaler (1986). A first primary statement of the broader combined expertise/experience/market question was not verified.","earliest_verified_reference":null,"current_reference":{"authors":["Benjamin Enke","Uri Gneezy","Brian Hall","David Martin","Vadim Nelidov","Theo Offerman","Jeroen van de Ven"],"title":"Cognitive Biases: Mistakes or Missing Stakes?","year":2023,"url":"https://benjamin-enke.com/pdf/Biases_stakes.pdf","doi":"10.1162/rest_a_01093","locator":"§I, p. 818; §IV Discussion, p. 830 (PDF pp. 1 and 13).","evidence_summary":"The paper states and tests the question whether larger stakes reduce four cognitive biases. Its discussion reports that the increase has small or no effects; it does not explicitly pose the combined expertise, experience, and actual-market extension."},"formalization_note":"Proposed empirical extension. The four biases in the cited experiment already have high-stakes evidence; their experimental result itself is not an open problem."},"statusQualification":"Provenance pending: an explicit published statement of this exact open question was not verified. This is a catalogue-authored candidate specification, not a certified open conjecture. Proposed empirical extension. The four biases in the cited experiment already have high-stakes evidence; their experimental result itself is not an open problem.","statusReviewedAt":"2026-10-08"}
% FIRST_POSTED: 2026-10-08
% ENTRY_KIND: statement_only
% !TeX program = lualatex
\documentclass[11pt]{article}
\usepackage{fontspec}
\usepackage[a4paper,margin=25mm]{geometry}
\usepackage{amsmath,amssymb,mathtools}
\usepackage{xurl}
\usepackage[hidelinks]{hyperref}
\newcommand{\sref}[2]{\hyperlink{source:#1:#2}{[S#2]}}
\setlength{\emergencystretch}{3em}
\begin{document}
% problemId:30007584
\section*{Biases under high stakes and experience}\label{30007584}
\subsection{Definitions and mathematical statement}
Fix a decision task \(k\) with objectively correct response \(y_k^*(X)\) for problem attributes \(X\). Let \(Y_{ik}(s,t)\) be individual \(i\)'s response when assigned incentive level \(s\geq0\) after \(t\) opportunities to practice with feedback. Define a preregistered bias score \(B_{ik}(s,t)=b_k(Y_{ik}(s,t),y_k^*(X),X)\), where \(b_k\) distinguishes a systematic directional error from unsystematic inaccuracy. Let \(H_i\) denote observed professional expertise and \(P\) a specified target population. The relevant response surface and incentive contrast are
\[\mu_{k,P}(s,t,h)=\mathbb E_P[B_{ik}(s,t)\mid H_i=h],\qquad \tau_{k,P}=\mu_{k,P}(s_H,t,h)-\mu_{k,P}(s_L,t,h).\]
Estimate these quantities using randomized incentives and feedback, adequate support for both stake levels, and payments expressed relative to participants' economic resources. Expertise is not automatically randomized; comparisons across \(h\) require an explicit selection or causal design. Which tasks retain materially positive bias at high stakes after meaningful experience, and how accurately do laboratory estimates predict bias in matched consequential market decisions? Specify thresholds for economic importance and distinguish learning from selective attrition, fatigue, and changes in task interpretation. Existing high-stakes experiments answer particular task-population contrasts; they do not identify this whole response surface or warrant a claim that every bias persists in every professional setting.

\par\textbf{Formulation and scope.} Proposed empirical extension. The four biases in the cited experiment already have high-stakes evidence; their experimental result itself is not an open problem.

\par\textbf{Open-question provenance.} An explicit published open-question statement was not verified. Sources below are supporting literature and must not be treated as a first listing.

\par\textbf{Historical priority.} The paper tests an older high-stakes question and cites Thaler (1986). A first primary statement of the broader combined expertise/experience/market question was not verified.

\subsection{Short English statement}
Estimate how stakes and repeated experience change the specified bias measures, distinguishing learning from selection and preference changes.

\subsection{Sources}
\begin{itemize}\raggedright
\item[\textnormal{[S1]}]\hypertarget{source:30007584:1}{} Benjamin Enke, Uri Gneezy, Brian Hall, David Martin, Vadim Nelidov, Theo Offerman, Jeroen van de Ven. 2023. Cognitive Biases: Mistakes or Missing Stakes?. §I, p. 818; §IV Discussion, p. 830 (PDF pp. 1 and 13)..
\url{https://benjamin-enke.com/pdf/Biases_stakes.pdf}.
DOI: 10.1162/rest\_a\_01093.
{\small\textit{Supporting or current literature; see evidence qualification. The paper states and tests the question whether larger stakes reduce four cognitive biases. Its discussion reports that the increase has small or no effects; it does not explicitly pose the combined expertise, experience, and actual-market extension.}}
\end{itemize}
\end{document}
